% from ega2-1-fr.tex
\oldpage[II]{117}

Comme le sous-$\mathcal{A}$-Module de $\mathcal{B}$ engendré par les $f^n$ est
une $\mathcal{A}$-Algèbre, il est clair que a) entraîne b). D'autre part, b)
implique que pour tout $x\in X$, le $\mathcal{A}_x$-module $\mathcal{C}_x$ est
de type fini, ce qui entraîne que tout élément de l'algèbre $\mathcal{C}_x$ et
en particulier $f_x$, est entier sur $\mathcal{A}_x$. Enfin, si pour un
$x\in X$, on a une relation de la forme
\[
f_x^n+(a_1)_xf_x^{n-1}+\ldots+(a_n)_x=0
\]
les $a_i$ ($1\leq i\leq n$) étant des sections de $\mathcal{A}$ au-dessus
d'un voisinage ouvert $U$ de $x$, la section
$f^n|U+a_1\cdot f^{n-1}|U+\ldots+a_n$ est nulle au-dessus d'un voisinage
$V\subset U$ de $x$, d'où résulte aussitôt que tous les $f^k|V$ ($k\geq0$)
sont combinaisons linéaires à coefficients dans $\Gamma(V,\mathcal{A})$ des
$f^j|V$ pour $0\leq j\leq n-1$~; on en conclut que c) entraîne a).

Lorsque les conditions équivalentes de (\hyperref[II.6.3.1-fr]{6.3.1}) sont
remplies on dit que la section $f$ est \emph{entière sur $\mathcal{A}$}.

\begin{corollary}[6.3.2]
\label{II.6.3.2-fr}
Sous les hypothèses de (\hyperref[II.6.3.1-fr]{6.3.1}), il existe un
sous-$\mathcal{A}$-Module (unique) $\mathcal{A}'$ de $\mathcal{B}$ tel que pour
tout $x\in X$, $\mathcal{A}'_x$ soit l'ensemble des germes
$f_x\in\mathcal{B}_x$ qui sont entiers sur $\mathcal{A}_x$. Pour tout ouvert
$U\subset X$, les sections de $\mathcal{A}'$ au-dessus de $U$ sont les
sections $f\in\Gamma(U,\mathcal{B})$ qui sont entières sur
$\mathcal{A}|U$.
\end{corollary}

L'existence de $\mathcal{A}'$ est immédiate, en prenant pour
$\Gamma(U,\mathcal{A}')$ l'ensemble des $f\in\Gamma(U,\mathcal{B})$ tels que
$f_x$ soit entier sur $\mathcal{A}_x$ pour tout $x\in U$. La seconde
assertion résulte aussitôt de (\hyperref[II.6.3.1-fr]{6.3.1}).

Il est clair que $\mathcal{A}'$ est une sous-$\mathcal{A}$-Algèbre de
$\mathcal{B}$~; on dit que c'est la \emph{fermeture intégrale de
$\mathcal{A}$ dans $\mathcal{B}$}.

\begin{env}[6.3.3]
\label{II.6.3.3-fr}
Soient $(X,\mathcal{A})$, $(Y,\mathcal{B})$ deux espaces annelés,
\[
g=(\psi,\theta):X\to Y
\]
un morphisme. Soit $\mathcal{C}$ (resp. $\mathcal{D}$) une
$\mathcal{A}$-Algèbre (resp. une $\mathcal{B}$-Algèbre) et soit
\[
u:\mathcal{D}\to\mathcal{C}
\]
un $g$-morphisme (\hyperref[0.4.4.1-fr]{0, 4.4.1}). Alors, si
$\mathcal{A}'$ (resp. $\mathcal{B}'$) est la fermeture intégrale de
$\mathcal{A}$ (resp. $\mathcal{B}$) dans $\mathcal{C}$ (resp.
$\mathcal{D}$), la \emph{restriction} de $u$ à $\mathcal{B}'$ est un
$g$-morphisme
\[
u':\mathcal{B}'\to\mathcal{A}'.
\]

En effet, si $j$ est l'injection canonique $\mathcal{B}'\to\mathcal{D}$, il
suffit de montrer que
\[
v=u^\sharp\circ g^*(j):g^*(\mathcal{B}')\to\mathcal{C}
\]
applique $g^*(\mathcal{B}')$ dans $\mathcal{A}'$. Or, un élément de
$(g^*(\mathcal{B}'))_x=\mathcal{B}'_{\psi(x)}\otimes_{\mathcal{B}_{\psi(x)}}
\mathcal{A}_x$ est entier sur $\mathcal{A}_x$ par définition de
$\mathcal{B}'$, donc il en est de même de son image par $v_x$, ce qui établit
notre assertion.
\end{env}

\begin{proposition}[6.3.4]
\label{II.6.3.4-fr}
Soient $X$ un préschéma, $\mathcal{A}$ une $\mathcal{O}_X$-Algèbre
quasi-cohérente. La fermeture intégrale $\mathcal{O}'_X$ de
$\mathcal{O}_X$ dans $\mathcal{A}$ est alors une $\mathcal{O}_X$-Algèbre
quasi-cohérente, et pour tout ouvert affine $U$ de $X$,
$\Gamma(U,\mathcal{O}'_X)$ est la fermeture intégrale de
$\Gamma(U,\mathcal{O}_X)$ dans $\Gamma(U,\mathcal{A})$.
\end{proposition}

On peut se borner au cas où $X=\operatorname{Spec}(B)$ est affine et
$\mathcal{A}=\widetilde{A}$, $A$ étant une

